%! Confidence Intervals for the Difference of Two Proportions — Unit 6, Lesson 6.8 — Exit Ticket (Answer Key)
\documentclass[10pt]{article}
\usepackage{apstats-article}
\usepackage{apstats-key}   % pulls in -boxes; do NOT also load -boxes
\newcommand{\TallMath}[1]{$\displaystyle #1\rule[-1.4em]{0pt}{3.2em}$}

\begin{document}

\pageheader{Unit 6, Lesson 6.8}{Exit Ticket --- Answer Key}
\namedateperiod

\vspace{0.15in}

A school district surveyed random samples of middle and high school students about daily
reading habits. Among 160 middle schoolers, 72 reported reading for pleasure at least
30 minutes per day. Among 140 high schoolers, 49 reported the same.

Construct and interpret a 95\% confidence interval for $p_M - p_H$, the difference in the
proportions of middle and high school students who read for pleasure at least 30 minutes per day.

\textit{State:} \ansline{We will construct a 95\% two-sample $z$-interval for $p_M - p_H$, where $p_M$ = true proportion of middle schoolers and $p_H$ = true proportion of high schoolers who read for pleasure $\ge30$ min/day.}

\textit{Plan --- Check Conditions:}\\
\textbf{Random:} \ansline{Both samples were randomly selected --- stated.}\\
\textbf{Large Counts:} $160(0.45)=72$, $160(0.55)=88$, $140(0.35)=49$, $140(0.65)=91$ --- \ans{all $\ge10$. \checkmark}\\
\textbf{10\%:} \ansline{160 $\le10\%$ of all middle schoolers; 140 $\le10\%$ of all high schoolers --- assumed.}

\textit{Do:}\\
$\hat p_M = 72/160 = 0.45$; \quad $\hat p_H = 49/140 = 0.35$; \quad $\hat p_M - \hat p_H = \ans{0.10}$\\
$SE = \sqrt{0.45(0.55)/160 + 0.35(0.65)/140} = \sqrt{0.001547 + 0.001625} = \sqrt{0.003172} \approx \ans{0.0563}$\\
$CI = 0.10 \pm 1.960(0.0563) = 0.10 \pm 0.1103 = (\ans{-0.010}, \ans{0.210})$

\textit{Conclude:} \ansline{We are 95\% confident that the true difference in the proportions of middle and high school students who read for pleasure at least 30 minutes per day ($p_M - p_H$) is between $-0.010$ and $0.210$. Since the interval contains 0, we cannot conclude there is a significant difference.}

\end{document}
